\section{K-009: a uniform finite bound for diagonal lifts}
\label{sec:uniform-lift-bound}

For $n\ge2$, set
\[
 r=3n-2,\qquad B_n=\frac{(3-2/n)^{(3n-2)/2}}{n}.
\]
This bound is not the proposed sharp value $2^{n-1}$.
It supplies a finite, explicit search limit, including stabilizer lifts
whose order exceeds the order of their determinant character.

\begin{theorem}[K-009]
Let $S\subseteq\{1,\ldots,n\}^3$ be any tensor support.  The order of
the component group of its diagonal stabilizer is at most $\lfloor B_n\rfloor$.
If the support admits nonnegative masses with every one-coordinate
margin equal to one, its determinant-product character is trivial on
the connected stabilizer.  Every root of unity in its character image
therefore has a diagonal lift of order at most $\lfloor B_n\rfloor$.
For a two-power root, a two-power lift can be chosen with order at most
$2^{\lfloor\log_2 B_n\rfloor}$.

Consequently, a primitive root of order $m>B_n$ is not a pole of the
Hilbert series of $\mathbb C[(\mathbb C^n)^{\otimes3}]^{SL_n^3}$,
graded by ordinary degree divided by $n$.
\end{theorem}

\begin{proof}
Write $N=3n$, and associate to an edge $(i,j,k)$ the integer column
$e_i+e_{n+j}+e_{2n+k}$.  Let $L_S\subseteq\mathbb Z^N$ be the
lattice generated by the support columns.  The diagonal stabilizer
is the diagonalizable group with character group $\mathbb Z^N/L_S$.
Thus the order in question is the order of its torsion subgroup.

Choose a maximal rationally independent subset of the support
columns, with lattice $L_B$.  The group $L_S/L_B$ is a subgroup of
the torsion of $\mathbb Z^N/L_B$, because the two lattices have the
same rational span.  It follows that
\[
 |\operatorname{tors}(\mathbb Z^N/L_S)|
 \le |\operatorname{tors}(\mathbb Z^N/L_B)|.
\]
Extend these independent columns to $r=3n-2$ independent columns
chosen from the complete tensor incidence matrix; this is possible
because the complete matrix has rank $r$.  Let $L_{B'}$ be the
resulting lattice.  The subgroup $L_{B'}/L_B$ is free, and it has
zero intersection with the torsion of $\mathbb Z^N/L_B$.
Indeed a rational relation between a new column and the old span
would contradict independence of the enlarged column set.
Therefore the old torsion injects into
$\operatorname{tors}(\mathbb Z^N/L_{B'})$.
It suffices to bound the latter order, denoted $d$.

View $B'$ as an $N\times r$ matrix.  Its left kernel consists of
the two differences between the sums of the three coordinate
blocks.  These differences have primitive exterior product:
their two-by-two minors are zero for a pair of rows in one block
and $\pm1$ for a pair in different blocks.  Complementary-minor
duality, or expansion after choosing a nonzero maximal minor,
shows that all nonzero $r\times r$ minors of $B'$ have absolute
value $d$; their greatest common divisor is the torsion order
by the integer Smith normal form.  There are $3n^2$ such minors.
Cauchy--Binet consequently gives
\[
             \det((B')^t B')=3n^2d^2.                 \tag{K-009a}
\]

Every column has three ones, so the trace of this Gram matrix is
$3r$.  If $d_v$ are the row sums, its Rayleigh quotient on the
all-one vector is $\sum_v d_v^2/r$.  Since $\sum_vd_v=3r$,
Cauchy--Schwarz gives a largest eigenvalue $\Lambda\ge3r/n$.
The other $r-1$ positive eigenvalues have fixed sum $3r-\Lambda$.
The arithmetic--geometric mean inequality bounds their product.
The function $u((3r-u)/(r-1))^{r-1}$ decreases for $u\ge3$;
therefore
\[
 \det((B')^tB')
 \le \frac{3r}{n}\left(\frac{r}{n}\right)^{r-1}
 =3(3-2/n)^{3n-2}.
\]
Together with (K-009a), this gives $d\le B_n$.

The determinant-product character is represented by
$\chi={\bf1}_{3n}$.  A margin-one mass vector expresses $\chi$
as a rational linear combination of support columns: rational
coefficients may be chosen because the defining linear system has
rational coefficients and a nonempty rational polyhedron.
Hence $\chi$ is torsion in $\mathbb Z^N/L_S$ and is trivial on
the connected diagonal stabilizer.  Splitting the finitely generated
character group decomposes that stabilizer as a torus times a finite
diagonalizable group of order at most $B_n$.  Restricting to the
two-primary factor gives the asserted two-power lift bound.

Finally let a nonvertex quotient point be fixed by the grading
element of primitive order $m$.  Its unique closed $SL_n^3$ orbit
has a product-linear stabilizer element with determinant-product
character equal to that grading element: apply a product-linear
representative of the grading action, then an $SL_n^3$ element
identifying the two closed orbits.  The semisimple part still fixes
the tensor and has the same character.  Diagonalize its factors.
Diagonal-torus semistability supplies the margin-one condition.
Thus $m\le B_n$ by the preceding argument.

For $m>B_n$ the grading fixed scheme is supported only at the
vertex.  Its coordinate ring is finite-dimensional.  The Koszul
identity in ALPOGE\_GCT P-F23, equations (P-F23b)--(P-F23c), writes
the Hilbert series times factors nonzero at the selected root as
an alternating sum of Hilbert series of finite modules supported
there.  Those modules have finite length, so their Hilbert series
are polynomials.  The Hilbert series is therefore regular at the
selected root.  This proves the last assertion.
\end{proof}

For $n=4$, $B_4=3125/128=24.4140625$.  Thus order sixteen is an
exhaustive upper limit for two-primary lifts, independently of the
small-support analysis in K-007.  The complete modulus-sixteen
certificates in K-008 then prove the sharper bound eight for
two-power determinant characters in that dimension.  For general
$n$, the gap between this bound and $2^{n-1}$ remains open here.

The lattice estimate and its proof are elementary consequences of
Smith normal form, complementary minors and Cauchy--Binet; this is
not a claim that those methods or their general determinant bounds
are new.  Its use as a certified lift cutoff in this problem is
the contribution recorded under K-009.
